\documentclass[11pt,a4paper]{article}
%TC:macro \original [ignore,ignore,ignore,ignore]
%TC:macro \hinweis [ignore]
%TC:macro \station [ignore,ignore,ignore]
%TC:macro \weiter [ignore]
%TC:macro \rahmenbild [ignore,ignore]
\usepackage[T1]{fontenc}
\usepackage[utf8]{inputenc}
\usepackage[ngerman]{babel}
\usepackage{lmodern,microtype,amsmath,amssymb}
\usepackage[margin=22mm,top=20mm,bottom=21mm,headheight=15pt]{geometry}
\usepackage{graphicx,float,booktabs,tabularx,array,caption,tikz}
\usepackage{xcolor,fancyhdr,enumitem,hyperref}
\graphicspath{{../pictures/eaai_palette/}{../pictures/}{../pictures/jes_final/}}
\definecolor{tured}{HTML}{B6302D}
\definecolor{ink}{HTML}{333333}
\definecolor{secondaryfocus}{HTML}{39738E}
\hypersetup{colorlinks=true,linkcolor=ink,urlcolor=ink,
  pdftitle={Dissertationsvortrag: illustrierte Erzaehlfassung mit Formeln},
  pdfauthor={Florian Rzepka},bookmarksnumbered=true}
\captionsetup{font=small,labelfont=bf,justification=raggedright,singlelinecheck=false}
\setlength{\parindent}{0pt}
\setlength{\parskip}{6pt plus 1pt}
\setlength{\intextsep}{9pt}
\setlength{\textfloatsep}{10pt}
\setlist[itemize]{leftmargin=*,itemsep=4pt,topsep=5pt}
\pagestyle{fancy}
\fancyhf{}
\fancyhead[L]{\small Florian Rzepka\enspace |\enspace Dissertationsvortrag}
\fancyhead[R]{\small Illustrierte Erz\"ahlfassung}
\fancyfoot[L]{\small Arbeitsfassung vom 7. September 2026}
\fancyfoot[R]{\small\thepage}
\renewcommand{\headrulewidth}{0.3pt}
\newcommand{\station}[3]{\clearpage\section{#1}\label{station:\thesection}
  {\small\textit{#2\enspace |\enspace #3}}\par\vspace{3pt}}
\newcommand{\weiter}[1]{\clearpage{\large\bfseries #1}\par\vspace{4pt}}
\newcommand{\original}[4]{\begin{figure}[H]\centering
  \includegraphics[width=0.98\linewidth,height=#3mm,keepaspectratio]{#2}
  \caption*{\textbf{Dissertation, Abbildung #1.} #4}\end{figure}}
\newcommand{\hinweis}[1]{\par\vspace{2pt}{\small\color{ink}\textit{Vorbereitung:} #1}\par}
% The focus outlines are document annotations. The original bitmap is unchanged.
\newcommand{\rahmenbild}[2]{\begin{figure}[H]\centering
  \begin{tikzpicture}
    \node[anchor=south west,inner sep=0] (framework) at (0,0)
      {\includegraphics[width=.94\linewidth]{bms_requirements.png}};
    \begin{scope}[x={(framework.south east)},y={(framework.north west)}]
      \ifnum#1=1
        \draw[tured,line width=1.6pt] (.448,.470)--(.311,.692)--(.500,.998)--(.692,.692)--(.552,.470);
        \draw[tured,line width=1.2pt] (.502,.754) ellipse[x radius=.047,y radius=.075];
        \draw[secondaryfocus,line width=1.2pt,dashed] (.808,.133) ellipse[x radius=.041,y radius=.057];
      \fi
      \ifnum#1=2
        \draw[tured,line width=1.6pt] (.663,.391)--(.806,.619)--(.998,.313)--(.806,.007)--(.663,.235);
        \draw[tured,line width=1.6pt] (.350,.390)--(.208,.616)--(.018,.310)--(.208,.004)--(.350,.230);
      \fi
      \ifnum#1=3
        \draw[tured,line width=1.6pt] (.448,.470)--(.311,.692)--(.500,.998)--(.692,.692)--(.552,.470);
        \draw[tured,line width=1.6pt] (.350,.390)--(.208,.616)--(.018,.310)--(.208,.004)--(.350,.230);
        \draw[secondaryfocus,line width=1.2pt,dashed] (.808,.133) ellipse[x radius=.041,y radius=.057];
      \fi
    \end{scope}
  \end{tikzpicture}
  \caption*{\textbf{Dissertation, Abbildung 2.1.} #2 Rot markiert den Hauptschwerpunkt.\ifnum#1=2\else\ Ein gestrichelter blauer Kreis markiert die erg\"anzend betrachtete Genauigkeit.\fi}
  \end{figure}}
\newcommand{\clip}{\operatorname{clip}}
\newcommand{\MAE}{\operatorname{MAE}}
\newcommand{\SOH}{\operatorname{SOH}}
\newcommand{\SOC}{\operatorname{SOC}}
\begin{document}
%TC:ignore
\thispagestyle{empty}
{\small\color{tured} DISPUTATION\enspace /\enspace INHALTLICHE ARBEITSFASSUNG}
\vspace{6mm}

{\LARGE\bfseries Neuronale Zustandsbestimmung\par
f\"ur eingebettete\par
Batteriemanagementsysteme\par}
\vspace{4mm}
{\large Illustrierte Erz\"ahlfassung mit Modellgrundlagen und Formeln}

\vspace{5mm}
\textbf{Leitgedanke.} Welche Information ben\"otigt ein neuronales Modell, wie zuverl\"assig verarbeitet es diese Information und wie wird daraus eine kompakte, praktisch einsetzbare BMS-Implementierung? Diese drei Fragen verbinden Datenrepr\"asentation, Robustheit und Embedded-Optimierung zu einer gemeinsamen Arbeit.

\textbf{Lesart.} Die Erz\"ahlung folgt 27 Bildstationen mit drei vorgeschalteten Schwerpunktansichten. Die Abbildungsnummern sind durchg\"angig die Nummern der \emph{Dissertation}. Die Bilder stammen aus der aktuellen Dissertation. Beim BMS-Anforderungsbild markieren zus\"atzliche Konturen den jeweiligen Schwerpunkt, ohne die Originaldatei zu ver\"andern. Die Stationen sind noch keine fertigen Folien.

\renewcommand{\arraystretch}{1.3}
\begin{tabularx}{\linewidth}{@{}>{\raggedright\arraybackslash}Xrr@{}}
\toprule
Abschnitt & Stationen & Zeit \\
\midrule
BMS und experimentelle Grundlage & 1 bis 3 & 3:30 \\
Zeitliche Repr\"asentation mit einem MLP & 4 bis 8 & 6:30 \\
Rekurrente Modelle und Robustheitsbenchmark & 9 bis 17 & 8:15 \\
Ausf\"uhrung auf dem Mikrocontroller & 18 bis 19 & 1:45 \\
Pruning, Quantisierung und gemessene Effekte & 20 bis 26 & 6:45 \\
Gesamtergebnis und Ausblick & 27 & 2:15 \\
\midrule
Erz\"ahlung & & 29:00 \\
Puffer & & 1:00 \\
\bottomrule
\end{tabularx}

\vspace{3mm}
\textbf{Roter Faden.} Modell gestalten, Verhalten pr\"ufen, Implementierung optimieren. Beim MLP liegt der Schwerpunkt auf Deployment und Modellkomplexit\"at. Der Benchmark verbindet Performance und Hardware. Die Optimierung verbindet Deployment und Hardware mit Accuracy als Kontrollgr\"o\ss e. Die Bereiche bleiben miteinander verbunden.

\textbf{Zeitliche Konsequenz.} Die Zeitbl\"ocke schlie\ss en die kurzen Schwerpunktwechsel ein. Pro Ergebnisbild wird eine Beobachtung erkl\"art, nicht jedes Teilpanel. Die Zeit ist ein Probenziel. Rauschvertiefung und vollst\"andige Gate-Gleichungen stehen als Reserve am Ende. Bei einer zu langen Generalprobe lassen sich die Syntheseansichten 6.15 und 6.16 gemeinsam besprechen.

\vfill
{\small Bildauswahl und fachliche Pr\"ufhinweise am Ende geh\"oren nicht zum gesprochenen Vortrag. Dissertation, Originalgrafiken und PowerPoint bleiben unver\"andert.}
%TC:endignore
% BEGIN SPOKEN NARRATIVE
\station{BMS-Anforderungen und Forschungsziel}{1:15 min}{Anforderungen}
\label{sec:start}
Wie l\"asst sich komplexes Batterieverhalten mit neuronalen Netzen beschreiben, die zugleich auf einem begrenzten Mikrocontroller betrieben werden k\"onnen? Diese Frage steht im Zentrum meiner Arbeit. Ein BMS muss den Ladezustand und den Alterungszustand aus verf\"ugbaren Messungen sch\"atzen. Beide Gr\"o\ss en beschreiben unterschiedliche Eigenschaften:
\begin{equation}
  \SOC(t)=\frac{Q_{\mathrm{available}}(t)}{C_{\mathrm{current}}(t)},
  \qquad
  \SOH(t)=\frac{C_{\mathrm{current}}(t)}{C_{\mathrm{nominal}}}.
\end{equation}
Eine gealterte Batterie kann voll geladen sein und trotzdem weniger Ladung speichern als eine neue. Gerade in dezentralen Speichern und mobilen Anwendungen muss diese Unterscheidung langfristig lokal verf\"ugbar bleiben.

\original{2.1}{bms_requirements.png}{91}{BMS-Anforderungen als gemeinsamer Bezugspunkt der Arbeit.}

Deployment beschreibt die Gestaltung und \"Ubertragung eines Modells in eine anpassbare BMS-Implementierung. Performance beschreibt die G\"ute seiner Zustandsbestimmung: Genauigkeit, Robustheit und Erholung. Hardware erfasst die daf\"ur ben\"otigten Ressourcen: Flash, RAM und Inferenzzeit. Modellkomplexit\"at ist somit eine Entwurfsentscheidung, Speicher und Laufzeit sind ihre konkreten Auswirkungen.

Diese Anforderungen bleiben in der gesamten Arbeit verbunden. Zuerst gestalten wir ein Modell mit geeigneter zeitlicher Information. Danach pr\"ufen wir sein Verhalten unter gest\"orten Bedingungen. Abschlie\ss end optimieren wir die Implementierung. Das Anforderungsbild zeigt jeweils, welcher dieser Aspekte gerade im Vordergrund steht.

\station{Daten als Grundlage der Zustandsbestimmung}{1:30 min}{Datenbasis}
Die Modelle ben\"otigen Messreihen, in denen nicht nur einzelne Betriebszust\"ande, sondern auch die Entwicklung der Batterie sichtbar wird. Die erste Untersuchung verwendet 14 NMC-Zellen unter unterschiedlichen Belastungsprofilen. Abbildung 4.1 zeigt die Alterung sowohl \"uber Zyklen als auch \"uber die verstrichene Zeit.
\begin{figure}[H]\centering
  \includegraphics[width=.48\linewidth,height=52mm,keepaspectratio]{paper1_soh_cycles.png}\hfill
  \includegraphics[width=.48\linewidth,height=52mm,keepaspectratio]{paper1_soh_time.png}
  \caption*{\textbf{Dissertation, Abbildung 4.1.} NMC-Alterungsverl\"aufe in beiden urspr\"unglichen Darstellungen.}
\end{figure}
F\"ur die folgenden Untersuchungen steht ein LFP-Datensatz mit 15 Zellen und acht Versuchsprofilen zur Verf\"ugung. Die Profile variieren unter anderem Lade- und Entladerate sowie Entladetiefe. Das Design of Experiments verteilt die gew\"ahlten Bedingungen systematisch, damit nicht nur ein einzelner Betriebsfall beobachtet wird.
\original{4.4}{Figure_26_APPENDIX_SOH_Aging_Conditions.png}{88}{LFP-Zellen, Alterungsverl\"aufe und unterschiedliche Betriebsbedingungen.}
Diese Vielfalt ist wichtig, weil ein Modell Alterung von kurzfristigen Last- und Temperatureffekten unterscheiden muss. F\"ur den Robustheitsbenchmark werden Trainings-, Validierungs- und Testzellen getrennt. Die sechs Testzellen pr\"ufen somit die \"Ubertragung auf andere Zellen, nicht nur auf sp\"atere Messpunkte derselben Trainingsreihe.

\station{Vom Messsystem zum programmierbaren BMS}{0:45 min}{Hardwarebasis}
Ein trainiertes Netz ist noch kein Batteriemanagementsystem. Dazu geh\"oren die Messung der Zellgr\"o\ss en, die \"Uberwachung und ein Controller, auf dem die Zustandsbestimmung ausgef\"uhrt werden kann. Die entwickelte Vier-Zell-Plattform verbindet einen Batterie\"uberwachungsbaustein mit einem programmierbaren STM32-Mikrocontroller.
\original{4.6}{bms_board_render.png}{110}{Eigene BMS-Plattform mit Mess- und Verarbeitungselektronik.}
Der TLE9012DQU \"ubernimmt die zellnahe \"Uberwachung. Der STM32H755 stellt die Umgebung f\"ur die eigene Software bereit. Die sp\"ater gezeigten Laufzeit- und Speichermessungen wurden auf einem STM32H753 durchgef\"uhrt. Die eigene Platine und dieser Benchmark-Controller haben damit unterschiedliche Rollen. Gemeinsam machen sie deutlich, wo die Methoden eingesetzt werden sollen: zwischen realen Messsignalen und einer begrenzten lokalen Recheneinheit.

\weiter{Zeitliche Repr\"asentation: Modellgestaltung}
\label{focus:mlp}
\rahmenbild{1}{Schwerpunkt Deployment, insbesondere Modellkomplexit\"at. Accuracy dient der Bewertung der gew\"ahlten Repr\"asentation.}
Der erste Teil setzt bei der Modellgestaltung an. Wie l\"asst sich zeitabh\"angiges Batterieverhalten bereits mit einer einfach aufgebauten neuronalen Architektur erfassen? Dazu untersuchen wir, welche Information wir dem Netz durch Merkmalsbildung und eine explizite Historie bereitstellen k\"onnen. Die erreichbare Genauigkeit zeigt anschlie\ss end, ob diese Repr\"asentation funktioniert.

Im Anforderungsbild steht daher die Modellkomplexit\"at im Vordergrund. Hardwareeffizienz ist das \"ubergeordnete Entwurfsziel. Eine konkrete MCU-Vermessung folgt in den weiteren Teilen der Arbeit.

\station{Ein MLP als Ausgangspunkt}{1:30 min einschlie\ss lich Schwerpunktwechsel}{Neuronale Grundlagen}
Zun\"achst stellt sich die Frage, wie viel Modellstruktur \"uberhaupt n\"otig ist. Ein Multilayer Perceptron ist ein geeigneter Ausgangspunkt, weil es aus einer Folge gut beschreibbarer Rechenoperationen besteht. Eine Schicht gewichtet ihre Eing\"ange, addiert einen Bias und wendet eine Aktivierungsfunktion an:
\begin{equation}
  \mathbf a^{(\ell)}=f^{(\ell)}\!\left(
  W^{(\ell)}\mathbf a^{(\ell-1)}+\mathbf b^{(\ell)}\right),
  \qquad \operatorname{ReLU}(z)=\max(0,z).
\end{equation}
\original{3.2}{mlp_general_architecture.png}{87}{Eing\"ange, verborgene Schichten und Ausgabe eines MLP.}
Die trainierbaren Parameter sind die Gewichte und Bias-Werte. F\"ur eine voll verbundene Schicht mit $n_{\mathrm{in}}$ Eing\"angen und $n_{\mathrm{out}}$ Ausg\"angen gilt:
\begin{equation}
  N_{\mathrm{Parameter}}=n_{\mathrm{in}}n_{\mathrm{out}}+n_{\mathrm{out}}.
\end{equation}
Mehr Neuronen bedeuten damit nicht nur mehr Darstellungsf\"ahigkeit, sondern auch mehr Speicher und Rechenaufwand. Ein MLP besitzt jedoch keinen internen zeitlichen Zustand. Werden nur aktuelle Messwerte eingegeben, steht die vorherige Belastung nicht automatisch zur Verf\"ugung. Deshalb betrachten wir zun\"achst die Eingangsrepr\"asentation, bevor wir die Architektur aufwendiger machen.

\station{Welche Information steckt in den Messgr\"o\ss en?}{1:00 min}{Merkmalsbildung}
Abbildung 5.1 liefert einen ersten Blick auf die Beziehungen im NMC-Datensatz. Die kumulierten Lade- und Entlademengen weisen eine starke negative Korrelation mit dem SOH auf. F\"ur beide liegt der Korrelationskoeffizient ungef\"ahr bei $-0{,}99$. F\"ur Temperatur liegt er bei $-0{,}58$, f\"ur Spannung in dieser Auswertung nahe null.
\original{5.1}{paper1_correlation.png}{95}{Korrelationsmatrix der Messgr\"o\ss en und abgeleiteten Merkmale.}
\begin{equation}
  r_{xy}=\frac{\sum_i(x_i-\bar x)(y_i-\bar y)}
  {\sqrt{\sum_i(x_i-\bar x)^2}\sqrt{\sum_i(y_i-\bar y)^2}}.
\end{equation}
Die Ladungsmengen enthalten bereits Belastungsgeschichte. Damit stellen sie dem Netz Information bereit, die im momentanen Stromwert allein fehlt. Diese Beobachtung motiviert Feature Engineering. Sie beweist allerdings weder eine kausale Wirkung noch die Wichtigkeit eines Eingangs im trainierten Netz. Insbesondere kann eine Gr\"o\ss e mit geringer linearer Korrelation in Kombination mit anderen Eing\"angen dennoch n\"utzlich sein.

\station{Zeitlicher Kontext ohne rekurrente Architektur}{1:00 min}{Lag-Sequenzen}
Wir erweitern den Eingang des MLP um eine geordnete Historie. Zuerst werden hoch aufgel\"oste Messpunkte zu gr\"o\ss eren Zeitintervallen zusammengefasst. Der Merkmalsvektor eines solchen Intervalls bildet einen Lag-Step. Mehrere aufeinanderfolgende Lag-Steps ergeben eine Lag-Sequenz.
\original{5.2}{paper1_lag_sequence.png}{94}{Zeitliche Aggregation und Aufbau der Lag-Sequenz.}
\begin{equation}
  \mathbf x_k=\begin{bmatrix}
  \boldsymbol\phi_{k-L+1}^{\mathsf T}&\cdots&
  \boldsymbol\phi_k^{\mathsf T}
  \end{bmatrix}^{\mathsf T},
  \qquad \widehat{\SOH}_k=f_{\theta}(\mathbf x_k).
\end{equation}
Hier bezeichnet $\boldsymbol\phi_k$ die Merkmale eines Zeitintervalls und $L$ die Anzahl dieser Intervalle. Beispielsweise bilden 16 Schritte mit jeweils zehn Minuten eine Historie von 160 Minuten. Ein einzelner aggregierter Punkt ist also noch keine Sequenz. Das MLP verarbeitet die gesamte Historie als einen gemeinsamen Eingangsvektor. Die zeitliche Information entsteht hier durch die Datenaufbereitung, nicht durch einen internen Speicher des Netzes.

\station{Training, Modellauswahl und unabh\"angige Evaluation}{1:30 min}{Methodischer Ablauf}
Wie daraus ein belastbarer Modellvergleich entsteht, zeigt der Ablauf in Abbildung 5.3. Die Messdaten werden bereinigt, Merkmale gebildet und zeitlich geordnet. Die Skalierung bringt Eing\"ange mit unterschiedlichen Einheiten in einen geeigneten Wertebereich. Ihre Parameter werden aus den Trainingsdaten bestimmt und anschlie\ss end unver\"andert auf weitere Daten angewendet.
\original{5.3}{paper1_process.png}{110}{Datenaufbereitung, Training, Hyperparameteroptimierung und Evaluation.}
Training passt die Gewichte an. Validierung dient der Auswahl von Einstellungen wie Netzbreite, Sequenzl\"ange und Zeitaufl\"osung. Die zur abschlie\ss enden Bewertung reservierten Testdaten d\"urfen diese Auswahl nicht steuern. Hyperparameteroptimierung untersucht daher mehrere Konfigurationen, statt nur ein Netz mit willk\"urlich gew\"ahlten Einstellungen zu trainieren.
\begin{equation}
  \mathcal L_{\mathrm{MSE}}=\frac1N\sum_{i=1}^{N}(y_i-\hat y_i)^2,
  \qquad
  \MAE=\frac1N\sum_{i=1}^{N}|y_i-\hat y_i|.
\end{equation}
Der quadratische Fehler kann das Training steuern. Der MAE beschreibt die mittlere absolute Abweichung f\"ur die Evaluation. Bei einem auf null bis eins normierten SOH entsprechen $0{,}011$ genau 1,10 Prozentpunkten. Die Trennung von Lernen, Auswahl und Pr\"ufung bildet auch f\"ur die nachfolgenden neuronalen Modelle die methodische Grundlage.

\station{Einfluss der Historie und erreichbare SOH-Sch\"atzung}{1:30 min}{MLP-Ergebnisse}
\original{5.5}{paper1_mae_matrix.png}{73}{Fehlervergleich f\"ur unterschiedliche Kombinationen aus Lag-L\"ange und Zeitaufl\"osung.}
\original{5.6}{paper1_soh_pred.png}{73}{Vorhergesagter SOH und Referenzverl\"aufe f\"ur ausgew\"ahlte Testf\"alle.}
Die beiden Bilder beantworten zusammen die Ausgangsfrage. Erstens beeinflusst die Repr\"asentation das Ergebnis. Eine l\"angere Historie ist nicht automatisch besser. In der untersuchten Suche ergab die Kombination aus 16 Schritten und zehn Minuten eine geeignete Konfiguration. Zweitens l\"asst sich der Alterungsverlauf bereits mit einer Feedforward-Architektur abbilden. Der mittlere Test-MAE liegt bei 1,10 Prozentpunkten, einzelne atypische Verl\"aufe sind schwieriger.

Damit ist gezeigt, dass zeitliche Merkmale schon ohne rekurrente Zellen n\"utzlich sind. Das ausgew\"ahlte MLP besitzt allerdings mehrere breite Schichten. Eine einfache \emph{Architekturfamilie} bedeutet nicht automatisch wenig Parameter oder nachgewiesene MCU-Eignung. Der n\"achste Schritt untersucht deshalb eine andere Form der zeitlichen Verarbeitung und anschlie\ss end deren konkrete Implementierungskosten.

\station{Von expliziten Lag-Sequenzen zu rekurrentem Kontext}{1:45 min, beide Seiten}{GRU und LSTM}
\label{sec:recurrent}
Beim MLP mussten wir die Historie explizit zusammenstellen. Rekurrente Netze stellen eine wiederverwendbare Struktur zur Verf\"ugung, in der zeitliche Information intern fortgeschrieben wird. Das vereinfacht die Modellierung einer fortlaufenden Verarbeitung. Es beseitigt jedoch weder Feature Engineering noch die Notwendigkeit, mit geeigneten Sequenzen zu trainieren.
\original{3.3}{lstm_cell_eaai_style.png}{91}{LSTM mit Zellzustand, Hidden State und steuernden Gates.}
Beim LSTM entscheidet ein Forget Gate, wie viel des bisherigen Zellzustands erhalten bleibt. Ein Input Gate gewichtet neue Information. Ein Output Gate steuert, was davon als Hidden State weitergegeben wird:
\begin{align}
  \mathbf c_k&=\mathbf f_k\odot\mathbf c_{k-1}
             +\mathbf i_k\odot\mathbf g_k,\\
  \mathbf h_k&=\mathbf o_k\odot\tanh(\mathbf c_k).
\end{align}
Die Gate-Werte werden aus dem aktuellen Eingang und dem vorherigen Hidden State gelernt. $\odot$ bezeichnet elementweise Multiplikation. Der interne Zustand ersetzt damit nicht die Messdaten, sondern verdichtet deren bisherige Verarbeitung. Entscheidend ist anschlie\ss end, wann dieser Zustand weitergegeben oder zur\"uckgesetzt wird.

\weiter{GRU: zeitlicher Zustand mit weniger Gate-Struktur}
\original{3.4}{gru_cell_eaai_style.png}{106}{GRU mit Reset Gate, Update Gate und einem fortgeschriebenen Hidden State.}
Eine GRU verwendet keinen zus\"atzlichen separaten Zellzustand. In der hier dargestellten Konvention mischt das Update Gate den bisherigen Zustand und einen neuen Kandidaten:
\begin{equation}
  \mathbf h_k=(1-\mathbf z_k)\odot\mathbf h_{k-1}
              +\mathbf z_k\odot\widetilde{\mathbf h}_k.
\end{equation}
Damit bieten GRU und LSTM zwei verwandte Wege, mehrere Messgr\"o\ss en und ihren zeitlichen Zusammenhang zu verarbeiten. Der Wechsel vom MLP bedeutet also nicht, dass der erste Ansatz ungeeignet war. Er erweitert die Umsetzungsm\"oglichkeiten. Ob dieser gelernte Kontext im Betrieb tats\"achlich hilfreich ist, zeigt erst der Vergleich unter realistischen St\"orungen. Und ob er kompakt bleibt, muss auf der Zielhardware gepr\"uft werden.
\hinweis{Die vollst\"andigen Gate-Gleichungen stehen in der Reserve. Im Vortrag an den beiden Bildern den Informationsfluss erkl\"aren, nicht jede Matrix vorlesen.}

\weiter{Robustheitsbenchmark: Verhalten und Ressourcen}
\label{focus:benchmark}
\rahmenbild{2}{Schwerpunkt Performance und Hardware. Beide Bereiche werden nacheinander experimentell untersucht.}
Nach der zeitlichen Repr\"asentation richtet sich der Blick auf die Zuverl\"assigkeit im Betrieb. Wie verhalten sich unterschiedliche Zustandsbestimmungsmethoden bei ver\"anderten Messbedingungen, und welche Ressourcen ben\"otigen sie auf dem Mikrocontroller?

Zuerst vergleichen wir Genauigkeit, Robustheit und Erholung. Anschlie\ss end messen wir Flash, RAM und Laufzeit. Im Anforderungsbild werden daher Performance und Hardware gemeinsam hervorgehoben. Der Vergleich von Fensterbetrieb und kontinuierlicher Zustandsfortschreibung ber\"uhrt zus\"atzlich Deployment, weil sich dabei die Ausf\"uhrungsweise des neuronalen Modells \"andert.

\station{Vier Wege zur SOC-Bestimmung im selben Benchmark}{1:30 min einschlie\ss lich Schwerpunktwechsel}{Modellvergleich}
\original{6.1}{Figure_02_Methodology_Overview.png}{93}{Gemeinsamer Auswertungsablauf und vier SOC-Modellzweige.}
Der Benchmark vergleicht unterschiedliche Rechenprinzipien. Direct Measurement, DM, integriert den Strom mit einer festen Kapazit\"at. Hybrid Direct Measurement, HDM, verwendet zus\"atzlich einen gesch\"atzten SOH zur Kapazit\"atsanpassung. F\"ur positive Laderichtung und Zeit in Sekunden l\"asst sich der Kern schreiben als
\begin{equation}
  Q_{c,k}=Q_{c,k-1}+\frac{I_k\Delta t_k}{3600},\qquad
  \widehat{\SOC}_k=\clip\!\left(1+\frac{Q_{c,k}}{C_{\mathrm{eff},k}},0,1\right).
\end{equation}
$Q_c$ wird am Vollladeanker zur\"uckgesetzt. Bei DM ist $C_{\mathrm{eff}}=C_{\mathrm{nominal}}$, bei HDM $C_{\mathrm{eff}}=C_{\mathrm{nominal}}\widehat{\SOH}$.

Das hybride Ersatzschaltbildmodell HECM korrigiert seine modellbasierte Vorhersage zus\"atzlich anhand der Spannung. Der datengetriebene Zweig DD verwendet eine GRU mit mehreren Eing\"angen, darunter Spannung, Strom, Temperatur, SOH und zeitliche Ableitungen. HDM, HECM und DD erhalten denselben LSTM-basierten SOH-Sch\"atzer. Unterschiede entstehen damit vor allem daraus, wie die SOC-Zweige die verf\"ugbare Information verarbeiten.

\station{Vom nominalen Test zu Betriebsst\"orungen}{0:30 min}{Benchmark-Design}
\original{6.2}{Figure_03_Disturbance_Taxonomy.png}{123}{St\"orungsfamilien im kausalen Benchmark.}
Die Testmatrix umfasst systematische Messfehler, zuf\"alliges Rauschen, Initialisierungsfehler, fehlende Daten und lokale Ereignisse. Alle Modelle erhalten vergleichbar gest\"orte Eing\"ange. Abgeleitete Gr\"o\ss en werden daraus neu berechnet. Die Auswertung folgt dem zeitlichen Verlauf und verwendet keine zuk\"unftigen Messwerte. So untersuchen wir nicht nur Genauigkeit unter idealen Eing\"angen, sondern auch Stabilit\"at und Erholung.

\station{Nominale Genauigkeit als Ausgangspunkt}{0:45 min}{Baseline}
\original{6.4}{Figure_04_Baseline_Performance.png}{112}{Baseline-Genauigkeit mit der Streuung zwischen den Testzellen.}
Zun\"achst wird ohne zus\"atzliche St\"orungen ausgewertet. Grundlage sind 16 ausgew\"ahlte Zeitfenster aus sechs Testzellen. Die Zellen gehen gleichgewichtet in das Gesamtergebnis ein. Dieser Zellmittelwert des MAE betr\"agt etwa 6,90 Prozentpunkte f\"ur DM, 4,12 f\"ur HDM, 3,37 f\"ur HECM und 2,58 f\"ur DD. Die Einzelpunkte zeigen die Streuung, die Fehlerbalken das 95-Prozent-Konfidenzintervall des aggregierten Werts.

DD erzielt hier die niedrigste beobachtete mittlere Abweichung. Das ist eine Eigenschaft der untersuchten Modelle und Testdaten, noch keine allgemeine Rangfolge aller Vertreter dieser Familien. Der interessante n\"achste Schritt ist deshalb, ob diese Reihenfolge unter ver\"anderten Messbedingungen erhalten bleibt und welche Fehlermechanismen dabei sichtbar werden.

\station{Warum systematische Stromfehler kritisch sind}{1:00 min}{Current-Gain-Sensitivit\"at}
\original{6.5}{Figure_05_Current_Bias_REVISED.png}{104}{Einfluss eines multiplikativen Strommessfehlers auf den SOC-Fehler.}
Ein Stromsensor kann einen Verst\"arkungsfehler, einen additiven Offset oder Rauschen aufweisen. Diese drei Anteile lassen sich unterscheiden durch
\begin{equation}
 I_{\mathrm{mess},k}=(1+g)I_k+b+\varepsilon_k.
\end{equation}
Abbildung 6.5 untersucht $g$, also den Gain. Bei drei Prozent steigt der MAE um etwa 1,08 Prozentpunkte f\"ur DM, 1,01 f\"ur HDM, 0,60 f\"ur HECM und 0,38 f\"ur DD. Ein systematischer Fehler im Strom geht bei der Ladungsbilanz fortlaufend in die aufsummierte Ladung ein.

HECM besitzt zus\"atzlich Spannungsr\"uckkopplung, bleibt aber vom Strom abh\"angig. DD verarbeitet weitere Messgr\"o\ss en und deren zeitlichen Kontext. Das bietet zus\"atzliche Information, ersetzt jedoch keinen fehlerfreien Sensor. Der additive Offset $b$ wird separat untersucht und erweist sich in der Gesamt\"ubersicht als besonders kritisch.

\station{Erholung durch zus\"atzliche Zustandsinformation}{1:00 min}{Initialisierungsfehler}
\original{6.9}{Figure_09_Initial_State_Recovery_Six_Cell.png}{98}{Erholung nach einer anf\"anglichen Zustandsabweichung.}
Eine falsche Initialisierung verschiebt den Startpunkt der Sch\"atzung. Reines Stromz\"ahlen kann diese Verschiebung nicht aus sich heraus erkennen. HECM verwendet dagegen die Differenz zwischen gemessener und vorhergesagter Spannung zur Korrektur:
\begin{equation}
  \mathbf x_k^{+}=\mathbf x_k^{-}
    +\mathbf K_k\left(U_{\mathrm{mess},k}-\widehat U_k^{-}\right).
\end{equation}
Die Zustandsgr\"o\ss en umfassen SOC und die beiden RC-Spannungen. Die Gewichtung $\mathbf K_k$ legt fest, wie stark die Spannungsabweichung die Modellvorhersage korrigiert. Diese zus\"atzliche R\"uckkopplung erkl\"art, warum HECM trotz einer falschen Anfangslage wieder zum Referenzverlauf finden kann.

Beim untersuchten Fehler von minus zehn Prozentpunkten erreicht HECM das vorgegebene Erholungsband relativ zum ungest\"orten Vergleich dauerhaft nach im Mittel etwa 1,20 Stunden, DD nach 2,60 Stunden. DD kann den Zustand aus seiner mehrkanaligen Historie neu erschlie\ss en. Die Stromz\"ahler ben\"otigen dagegen einen geeigneten Anker. HECM zeigt hier also eine besondere St\"arke, obwohl DD den kleineren nominalen MAE besitzt.

\station{Lokale Ereignisse trotz kleinem Gesamtfehler}{0:45 min}{Spannungsspitze}
\original{6.13}{Figure_13_Voltage_Spike_Response_JES2.png}{117}{Spannungsspitze mit zeitaufgel\"oster Reaktion und ereignisbezogener Auswertung.}
Eine kurze Spannungsspitze liefert ein anderes Bild. DD kann unmittelbar stark reagieren, obwohl der zus\"atzliche MAE \"uber die gesamte Sequenz klein bleibt. Das ist mit der direkten Verarbeitung von Spannung und abgeleiteten dynamischen Merkmalen vereinbar. Der gelernte Kontext sch\"utzt also nicht automatisch vor jeder lokalen Eingabest\"orung.

DM und HDM verwenden die Spannung nicht als laufenden Korrektureingang. Ihre geringe Reaktion auf diesen Spannungsimpuls ist deshalb nicht einfach ein Mittelungseffekt der Stromintegration. Bei HECM wird die Spannungsinformation gewichtet in die Zustandskorrektur eingebracht. Das Beispiel zeigt, warum neben dem Gesamtfehler auch die unmittelbare Ereignisreaktion betrachtet werden muss.

\station{Welche St\"orungen bestimmen das Gesamtbild?}{0:30 min}{Szenario\"ubersicht}
\original{6.15}{Figure_14_Cross_Scenario_Heatmap_Six_Cell.png}{129}{Fehlerzunahme \"uber die untersuchten St\"orungsfamilien.}
Die \"Ubersicht ordnet die Einzelbeispiele ein. Vor allem systematische Stromfehler verursachen starke globale Abweichungen. Der additive Stromoffset f\"allt besonders auf. Kurze lokale Ereignisse k\"onnen dagegen im globalen Mittel klein erscheinen und trotzdem eine ausgepr\"agte kurzzeitige Reaktion ausl\"osen. Die beiden Betrachtungen beantworten unterschiedliche praktische Fragen und erg\"anzen sich.

\station{St\"arken der Modellfamilien statt einer einzigen Rangzahl}{0:30 min}{Benchmark-Synthese}
\original{6.16}{Figure_15_Decision_Synthesis_Six_Cell.png}{116}{Synthese von Accuracy, Robustness und Recovery.}
In der zusammengef\"uhrten Bewertung erzielt DD das st\"arkste Gesamtergebnis der untersuchten Vertreter. Die vorherigen Beispiele zeigen aber, warum die einzelnen Dimensionen erhalten bleiben m\"ussen. HECM kann bei Erholung besonders gut abschneiden, w\"ahrend DD bei lokalen Eingangsspitzen empfindlich sein kann. Die normierten Scores im Bild sind relative Vergleichsgr\"o\ss en, keine Genauigkeitsprozente. F\"ur das BMS folgt daraus, ein ausgew\"ahltes Modell gezielt auf seine relevanten Schwachstellen zu pr\"ufen und darauf vorzubereiten.

\station{Die Methoden auf dem Mikrocontroller}{0:45 min}{Embedded-Performance}
\original{6.19}{Figure_19_Hardware_Memory_Footprints.png}{95}{Flash und zur Laufzeit bestimmter RAM-Bedarf der isolierten SOC-Implementierungen.}
Nun wechseln wir im Anforderungsrahmen von der Sch\"atzg\"ute zur Ausf\"uhrung. Alle vier SOC-Implementierungen wurden auf dem STM32H753 ausgef\"uhrt und mit ihren Software-Referenzen verglichen. Die Abweichungen sind sehr klein. Zugleich unterscheiden sich ihre Kosten deutlich.

DM und HDM ben\"otigen ungef\"ahr 27 KiB Flash, HECM rund 470 KiB und DD im kontinuierlichen Betrieb etwa 116 KiB. Der zugeh\"orige Peak-RAM liegt zwischen 2,4 und 4,1 KiB. Alle passen auf den gew\"ahlten Controller. Gemessen werden hier die SOC-Kerne. Der ben\"otigte SOH wird zugef\"uhrt, nicht mitsamt einem zus\"atzlichen SOH-Netz in diesen Zahlen versteckt.

\station{Fensterbetrieb und kontinuierliche Zustandsfortschreibung}{1:00 min}{Ausf\"uhrungsmodus}
\original{6.20}{Figure_20_DD_Inference_Mode_Comparison.png}{108}{DD-Vergleich von Rolling Window, Continuous State und Periodic Reset.}
Im Rolling-Window-Betrieb wird f\"ur jede neue Sch\"atzung erneut eine vollst\"andige Eingangshistorie verarbeitet. Hier umfasst sie ungef\"ahr 34 Minuten. Bei kontinuierlicher Ausf\"uhrung wird dagegen nur der neue Schritt berechnet und der bisherige rekurrente Zustand weitergegeben:
\begin{equation}
  \mathbf h_k=F_{\theta}(\mathbf x_k,\mathbf h_{k-1}),
  \qquad \hat y_k=G_{\theta}(\mathbf h_k).
\end{equation}
Das reduziert die gemessene DD-Laufzeit von etwa 724 Millisekunden auf 0,425 Millisekunden pro Ausgabe. Der Peak-RAM sinkt von rund 67,3 auf 4,1 KiB. Periodische Resets ben\"otigen ebenfalls wenig Ressourcen, erzeugen jedoch wiederkehrende Anlauftransienten.

Kontinuierlicher Betrieb ist damit ein sehr wirksamer Hebel. Er ist aber nicht mathematisch identisch mit dem wiederholten Fensterlauf, weil die Zustandsvorgeschichte anders ist. Die gute nominale \"Ubereinstimmung muss daher separat gepr\"uft werden. Die Robustheitsaussagen des urspr\"unglichen Fensterbenchmarks gelten nicht ungepr\"uft f\"ur den anderen Ausf\"uhrungsmodus.

\weiter{Embedded-Optimierung: Modell und Implementierung}
\label{focus:optimization}
\rahmenbild{3}{Schwerpunkt Deployment und Hardware. Die Genauigkeit bleibt als Kontrollgr\"o\ss e sichtbar.}
Der abschlie\ss ende Teil ver\"andert nun gezielt das Netz und seine Darstellung. Wie l\"asst sich eine funktionierende neuronale Zustandsbestimmung an begrenzte Hardwareressourcen anpassen, und welche Auswirkungen hat diese Optimierung auf die Sch\"atzg\"ute?

Pruning reduziert die Struktur, Quantisierung die Darstellungsbreite ausgew\"ahlter Gewichte. Das sind Deployment-Entscheidungen. Ihre Wirkung messen wir anhand von Flash, RAM und Laufzeit. Accuracy bleibt mit im Bild, weil eine Ressourceneinsparung nur gemeinsam mit ihrer Auswirkung auf die Zustandsbestimmung beurteilt werden kann.

\station{Welche Teile des Netzes lassen sich verkleinern?}{1:00 min einschlie\ss lich Schwerpunktwechsel}{LSTM-Architektur}
\original{7.1}{embedded_lstm_schematic.png}{82}{LSTM mit nachgeschaltetem MLP als Ausgangspunkt der Kompression.}
Der abschlie\ss ende Teil untersucht nun die Architektur selbst. Hier werden eigenst\"andige LSTM-Modelle f\"ur SOC und SOH verwendet, nicht unver\"andert die Modelle des vorherigen Benchmarks. Beide verarbeiten sechs Merkmale. Das SOC-Modell besitzt 64 rekurrente Einheiten, das SOH-Modell 128. Dahinter folgt ein MLP mit ebenfalls 64 beziehungsweise 128 Einheiten.
%TC:ignore
\begin{center}\small
\begin{tabular}{@{}lcc@{}}\toprule
Architektur und Pruning & SOC & SOH \\
\midrule
Eingangsdimension $D$ & 6 & 6 \\
LSTM-Hidden-Size $H$, Base & 64 & 128 \\
LSTM-Hidden-Size, Pruned & 45 & 90 \\
MLP-Breite $M$ & 64 & 128 \\
Ausgabe & SOC & SOH \\
\bottomrule\end{tabular}
\end{center}
%TC:endignore
Die Gewichte verbinden den aktuellen Eingang und den vorherigen Hidden State mit den Gates. Deshalb beeinflusst die rekurrente Breite mehrere Matrizen gleichzeitig. Hinzu kommen die Zust\"ande im RAM und die Anbindung an das MLP. Diese Struktur erkl\"art, warum das Entfernen ganzer rekurrenter Einheiten sowohl Speicher als auch Rechenzeit verringern kann.

\station{Von trainierten Parametern zu gepr\"ufter Firmware}{0:30 min}{Optimierungsablauf}
\original{7.3}{embedded_pipeline.png}{125}{Training, alternative Kompressionspfade, C-Export und Hardware-Evaluation.}
Aus dem trainierten FP32-Netz entstehen zwei getrennte Varianten. Beim strukturierten Pruning werden Einheiten entfernt und die verbleibenden Gewichte kurz nachtrainiert. Bei der Quantisierung werden ausgew\"ahlte Gewichtsmatrizen kompakter gespeichert. Anschlie\ss end werden die Parameter in C exportiert und auf derselben Hardware gepr\"uft. So lassen sich Auswirkungen auf Genauigkeit, Speicher und Laufzeit gemeinsam betrachten, statt nur eine kleinere Modelldatei zu melden.

\station{Strukturiertes Pruning und sein Rechenvorteil}{1:30 min}{Kompressionsgrundlagen}
\original{3.5}{pruning_granularity_comparison.png}{84}{Unstrukturiertes Nullsetzen im Vergleich zum Entfernen ganzer Strukturen.}
Einzelne Gewichte auf null zu setzen macht eine dichte Matrix noch nicht kleiner. Ohne passende Sparse-Kernel werden vielfach weiterhin dieselben Positionen verarbeitet. Strukturiertes Pruning entfernt dagegen ganze Einheiten samt ihren Verbindungen. Beim LSTM m\"ussen die zugeh\"origen Gate-Zeilen, rekurrenten Spalten und MLP-Eing\"ange konsistent verkleinert werden.

Die Auswahl basiert hier auf der Gewichtsnorm der rekurrenten Kan\"ale. Anschlie\ss end wird kurz nachtrainiert. Die Hidden Size sinkt um ungef\"ahr 30 Prozent. Der Einfluss auf die Zahl der Multiply-Accumulate-Operationen, kurz MACs, l\"asst sich absch\"atzen:
\begin{equation}
 N_{\mathrm{MAC}}\approx4H(D+H)+HM+M.
\end{equation}
Vier affine LSTM-Transformationen erkl\"aren den ersten Term. Die folgenden Terme geh\"oren zum MLP und zur Ausgabe. Da $H$ auch quadratisch eingeht, kann der Rechenaufwand st\"arker sinken als die Hidden Size. F\"ur SOC reduziert sich die Zahl von 22\,080 auf 12\,124 MACs pro Schritt, also um etwa 45 Prozent. Ob das auch schneller wird, pr\"ufen wir anschlie\ss end mit realer Laufzeitmessung.

\station{Quantisierung: weniger Bits, nicht automatisch weniger Laufzeit}{1:15 min}{Kompressionsgrundlagen}
\original{3.7}{quantization_scope_and_grids.png}{98}{Quantisierungsumfang und Abbildung auf ein diskretes Werteraster.}
Quantisierung ver\"andert die Darstellung der Gewichte. Anstelle eines FP32-Werts werden ein ganzzahliger Wert und ein Skalierungsfaktor verwendet. Hier wird jede Zeile der rekurrenten Gewichtsmatrizen symmetrisch auf INT8 abgebildet:
\begin{align}
 s_r&=\max\!\left(\frac{\max_j|w_{rj}|}{127},\epsilon\right),\\
 q_{rj}&=\clip\!\left(\operatorname{round}\frac{w_{rj}}{s_r},-127,127\right),
 \qquad \hat w_{rj}=s_rq_{rj}.
\end{align}
F\"ur diese gespeicherten Gewichte sinkt der Bedarf von vier auf ein Byte pro Wert, erg\"anzt um die Skalen. Das gesamte Netz wird dadurch aber nicht viermal kleiner. Zust\"ande, Aktivierungen, Bias und MLP bleiben in diesem Ansatz FP32.

Auch die Rechnung bleibt gleitkommabasiert. Die kompakten Gewichte werden w\"ahrend der Berechnung rekonstruiert. Das spart Flash, kann aber zus\"atzliche Laufzeit kosten. Pruning reduziert also die Struktur, Quantisierung die Darstellungsbreite. Beides sind unterschiedliche Hebel. Ihre Kombination ist eine sinnvolle Weiterentwicklung, wurde hier aber nicht als dritte kombinierte Kompressionsvariante vermessen.

\station{Wie viel Sch\"atzg\"ute bleibt erhalten?}{1:00 min}{SOC- und SOH-Ergebnisse}
\original{7.8}{embedded_soc_error.png}{66}{SOC-Fehlerverteilungen der drei Deployment-Varianten.}
\original{7.9}{embedded_soh_error.png}{66}{SOH-Fehlerverteilungen nach der jeweiligen kausalen Ausgabeverarbeitung.}
Bei SOC bleiben die Verteilungen in einer \"ahnlichen Gr\"o\ss enordnung. Die mittleren absoluten Fehler liegen bei ungef\"ahr 2,7 Prozentpunkten f\"ur Base, 2,3 f\"ur Pruned und 2,8 f\"ur Quantized. Pruning verschlechtert das Ergebnis hier also nicht zwangsl\"aufig. Das ist ein beobachtetes Ergebnis dieses Vergleichs, kein allgemeines Versprechen.

Bei SOH ist der Unterschied deutlicher. Base erreicht 0,85 Prozentpunkte, Pruned 1,46 und Quantized 1,41. Die komprimierten Modelle bleiben nutzbar, bezahlen ihre Ressourceneinsparung hier aber mit h\"oherem Fehler. Die Werte beziehen sich auf die beschriebene kausale Nachverarbeitung. Gerade bei einem langsamen SOH-Signal m\"ussen Gl\"attung und Reaktionsverz\"ogerung mitgedacht werden. Die Aussage lautet deshalb nicht, dass alle Varianten identisch sind, sondern dass sich unterschiedliche praktische Betriebspunkte erreichen lassen.

\station{Parameter, Flash und RAM gemeinsam betrachten}{0:45 min}{Speicherbedarf}
\original{7.12}{embedded_model_sizes.png}{131}{Parameterzahl, idealisierte Gewichtsspeicherung und Firmware-Ressourcen. Fachliche Pr\"ufhinweise zur Originalgrafik stehen am Ende dieser Fassung.}
Die Speicherauswertung trennt die reine Gewichtsdarstellung von der tats\"achlichen Firmware. Pruning senkt den SOC-Flash von rund 105 auf 62 KB und den SOH-Flash von 335 auf 182 KB. Quantisierung erreicht mit ungef\"ahr 52 und 138 KB die kleinsten Flash-Werte.

RAM umfasst dagegen auch Zust\"ande, Puffer und den Stack. Deshalb folgt seine Reduktion nicht einfach derselben Quote wie die Gewichtsspeicherung. Die wesentliche Botschaft ist, dass ein Modell nicht nur anhand seiner Parameterzahl beurteilt werden kann. Entscheidend bleibt die gemessene oder aus der Firmware bestimmte Ressource auf der konkreten Zielplattform.

\station{Weniger Rechenoperationen und gemessene Laufzeit}{0:45 min}{Inferenz und Gesamtlatenz}
\original{7.13}{embedded_latency_hist.png}{79}{Verteilung der End-to-End-Latenz mit UART-Kommunikation und Hardware-Ausf\"uhrung.}
Abbildung 7.13 zeigt die Zeit vom Host-Aufruf bis zur R\"uckgabe des Ergebnisses. Darin steckt auch Kommunikation. F\"ur die Optimierung des Netzes ist zus\"atzlich die direkt auf dem Controller gemessene Kernelzeit entscheidend:
%TC:ignore
\begin{center}\small
\begin{tabular}{@{}lrrr@{}}\toprule
Kernelzeit, Tabellen 7.1 und 7.2 & Base & Pruned & Quantized \\
\midrule
SOC [ms] & 1,40 & 0,80 & 6,99 \\
SOH [ms] & 22,73 & 12,72 & 29,21 \\
\bottomrule\end{tabular}
\end{center}
%TC:endignore
Pruning verk\"urzt die Netzberechnung bei beiden Aufgaben um ungef\"ahr 43 bis 44 Prozent. Quantisierung spart am meisten Flash, ben\"otigt in diesen gemischten INT8/FP32-Kernels aber l\"anger. Welcher Betriebspunkt geeigneter ist, h\"angt davon ab, ob Speicher, Zeit oder Genauigkeit die engste Grenze setzt. Das l\"asst sich aus einem PC-MAE allein nicht entscheiden.

\station{Gesamtergebnis und n\"achste Schritte}{2:15 min}{BMS-Anforderungen zusammenf\"uhren}
\label{sec:schluss}
Die Ergebnisse verbinden die drei Bereiche des Anforderungsbilds. Deployment beginnt bereits bei der Modellgestaltung. Das MLP-Beispiel zeigt, dass Merkmalsbildung und explizite Historie einen gro\ss en Beitrag leisten. Zeitabh\"angiges Batterieverhalten verlangt nicht automatisch die aufwendigste Netzstruktur. Rekurrente Modelle erweitern diesen Ansatz um einen intern fortgeschriebenen Zustand.

Performance beschreibt, wie verl\"asslich diese Modelle unter ver\"anderten Messbedingungen arbeiten. Im Benchmark bietet der datengetriebene Ansatz ein starkes Gesamtergebnis. Das passt zur gemeinsamen Verarbeitung mehrerer Eingangsinformationen und ihres zeitlichen Kontexts. Zugleich bleiben spezifische Schwachstellen sichtbar. Systematische Stromfehler belasten die Ladungsbilanz. Spannungsr\"uckkopplung hilft bei Erholung. Lokale Spannungsspitzen k\"onnen gerade im neuronalen Modell ausgepr\"agte Reaktionen hervorrufen. Diese Befunde liefern Ansatzpunkte f\"ur geeignete Trainingsst\"orungen, Signalbehandlung und Zustandsinitialisierung.

Auf der Hardware werden die Folgen der Entwurfsentscheidungen messbar. Kontinuierliche Zustandsfortschreibung vermeidet die wiederholte Berechnung einer langen Historie. Strukturiertes Pruning reduziert die Dimension des Netzes und damit Speicher und Laufzeit. Quantisierung verdichtet die Gewichtsdarstellung, ihre Geschwindigkeitswirkung h\"angt jedoch von der Rechenimplementierung ab. Diese Ma\ss nahmen m\"ussen daher zusammen mit ihrer Auswirkung auf die Sch\"atzung bewertet werden.

\original{2.1}{bms_requirements.png}{73}{R\"uckkehr zum Anforderungsrahmen: Information, Modellverhalten und Implementierung treffen auf demselben BMS zusammen.}

Der Gesamtbeitrag ist damit ein deploymentorientierter Weg vom Messsignal bis zur kompakten neuronalen Firmware. Die Netze sollen genug Komplexit\"at besitzen, um das Batterieverhalten zu erfassen, aber keine unn\"otigen Ressourcen beanspruchen. Die Hardware-Ergebnisse zeigen, dass dies mit den untersuchten Modellen auf einem begrenzten Mikrocontroller erreichbar ist.

Als n\"achste Schritte bieten sich die gemeinsame Optimierung von Architektur und Kompression sowie Hardware-in-the-Loop-Tests mit gezielt gest\"orten Messsignalen auf der eigenen BMS-Plattform an. Hinzu kommen breitere Betriebsdaten und eine erneute Robustheitspr\"ufung der kontinuierlichen Ausf\"uhrung. Damit wird die lokale Zustandsbestimmung nicht nur kompakter, sondern gezielter auf die Bedingungen vorbereitet, unter denen sie sp\"ater verl\"asslich arbeiten soll.

% END SPOKEN NARRATIVE
%TC:ignore
\clearpage
\section*{Reserve: Warum Stromrauschen anders wirkt}
\addcontentsline{toc}{section}{Reserve: Stromrauschen}
\original{6.7}{Figure_07_Noise_Robustness.png}{114}{Lokale Stromrauschreaktion und globaler Fehlervergleich. Reservebild zum Hauptbeispiel 6.13.}

Stromrauschen und eine Spannungsspitze illustrieren unterschiedliche Mechanismen. Bei DM und HDM geht jeder verrauschte Stromwert in die Ladungsbilanz ein. Positive und negative Fehler k\"onnen sich dabei teilweise ausgleichen. Es ist also nicht so, dass kleine neue Werte ignoriert werden. Ihr Beitrag wird integriert:
\begin{equation}
 \Delta Q_{c,N}=\frac{\Delta t}{3600}\sum_{k=1}^{N}\varepsilon_k.
\end{equation}
Bei unabh\"angigem, mittelwertfreiem Rauschen mit Varianz $\sigma_I^2$ gilt innerhalb eines Intervalls ohne Reset oder Begrenzung
\begin{equation}
 \mathbb E[\Delta Q_{c,N}]=0,\qquad
 \operatorname{Var}(\Delta Q_{c,N})=
 N\sigma_I^2\left(\frac{\Delta t}{3600}\right)^2.
\end{equation}
Die Standardabweichung w\"achst dann mit $\sqrt N$. Ein konstanter Offset liefert dagegen einen gerichteten Beitrag $Nb\Delta t/3600$. Zufallsrauschen verschwindet somit nicht vollst\"andig, erzeugt aber einen anderen Fehlerverlauf als ein anhaltender Offset.

DD verarbeitet Strom und Stromableitung direkt. Dadurch kann mehr kurzfristige Variation im Ausgang sichtbar sein, ohne dass der globale MAE stark steigt. HECM nutzt Strom auch in der Zustands- und Spannungsvorhersage. Ein Kalmanfilter ist daher nicht grunds\"atzlich gegen jede Form von Messrauschen unempfindlich.

\textbf{Auswahl f\"ur 30 Minuten.} Abbildung 6.13 bleibt im Hauptteil, weil sie den Unterschied zwischen globaler Genauigkeit und lokalem Ereignis besonders deutlich macht. Abbildung 6.7 ist sinnvoll als Reserve oder als Ersatz, wenn die Stromintegration im Vortrag st\"arker erkl\"art werden soll. Beide vollst\"andig zu pr\"asentieren w\"urde zus\"atzliche Zeit ben\"otigen.

\clearpage
\section*{Reserve: Gate-Gleichungen und Parameter}
\addcontentsline{toc}{section}{Reserve: Gate-Gleichungen}
\textbf{LSTM.} Aktueller Eingang $\mathbf x_k$, vorheriger Hidden State $\mathbf h_{k-1}$ und Zellzustand $\mathbf c_{k-1}$ bestimmen den n\"achsten Schritt:
\begin{align}
 \mathbf f_k&=\sigma(W_f\mathbf x_k+U_f\mathbf h_{k-1}+\mathbf b_f),\\
 \mathbf i_k&=\sigma(W_i\mathbf x_k+U_i\mathbf h_{k-1}+\mathbf b_i),\\
 \mathbf g_k&=\tanh(W_g\mathbf x_k+U_g\mathbf h_{k-1}+\mathbf b_g),\\
 \mathbf o_k&=\sigma(W_o\mathbf x_k+U_o\mathbf h_{k-1}+\mathbf b_o),\\
 \mathbf c_k&=\mathbf f_k\odot\mathbf c_{k-1}+\mathbf i_k\odot\mathbf g_k,\\
 \mathbf h_k&=\mathbf o_k\odot\tanh(\mathbf c_k).
\end{align}
Die Sigmoidfunktion $\sigma(z)=1/(1+e^{-z})$ beschr\"ankt die Gates auf null bis eins. Der Kandidat $\mathbf g_k$ ist neue Information, kein viertes Sigmoid-Gate.

\textbf{GRU.} Passend zur Darstellung in Abbildung 3.4 wird hier folgende Reset-before-Konvention verwendet:
\begin{align}
 \mathbf z_k&=\sigma(W_z\mathbf x_k+U_z\mathbf h_{k-1}+\mathbf b_z),\\
 \mathbf r_k&=\sigma(W_r\mathbf x_k+U_r\mathbf h_{k-1}+\mathbf b_r),\\
 \widetilde{\mathbf h}_k&=\tanh\!\left(W_h\mathbf x_k+
             U_h(\mathbf r_k\odot\mathbf h_{k-1})+\mathbf b_h\right),\\
 \mathbf h_k&=(1-\mathbf z_k)\odot\mathbf h_{k-1}
             +\mathbf z_k\odot\widetilde{\mathbf h}_k.
\end{align}
Softwarebibliotheken k\"onnen das Update Gate umgekehrt definieren oder den Reset nach der rekurrenten Matrixoperation anwenden. Beim C-Export ist deshalb die konkrete Operatorfolge des trainierten Modells ma\ss geblich. Das vereinfachte Zellbild ersetzt keine Implementierungsspezifikation.

\textbf{Speicherwirkung.} Bei Eingangsdimension $D$, Hidden Size $H$ und einem zusammengefassten Bias pro Transformation sind die Parameterzahlen der rekurrenten Schicht
\begin{equation}
 N_{\mathrm{LSTM}}=4H(D+H+1),\qquad
 N_{\mathrm{GRU}}=3H(D+H+1).
\end{equation}
Bei zwei getrennt gespeicherten Bias-Vektoren pro Transformation wird aus $+1$ jeweils $+2$. Hinzu kommen Ausgabe- und gegebenenfalls weitere Schichten. Die reinen FP32-Zust\"ande beanspruchen $8H$ Byte beim LSTM und $4H$ Byte bei der GRU. Das ist kein vollst\"andiger RAM-Wert, da Puffer, Stack und weitere Firmware hinzukommen.

\clearpage
\section*{Reserve: Repr\"asentation und Optimierung}
\addcontentsline{toc}{section}{Reserve: Repr\"asentation und Optimierung}
\textbf{Belastungshistorie als Merkmal.} Bei positiver Laderichtung und Zeitintervallen in Sekunden sind kumulierte Lade- und Entlademengen in Ah
\begin{align}
 Q_{+,k}&=Q_{+,k-1}+\frac{\max(I_k,0)\Delta t_k}{3600},\\
 Q_{-,k}&=Q_{-,k-1}+\frac{\max(-I_k,0)\Delta t_k}{3600}.
\end{align}
Anders als die vorzeichenbehaftete Ladungsbilanz enthalten beide Gr\"o\ss en den jeweiligen akkumulierten Durchsatz. Die hohe Korrelation mit Alterung motiviert ihren Einsatz, kann aber zugleich durch die gemeinsame Entwicklung mit der Versuchsdauer gepr\"agt sein.

\textbf{Skalierung.} F\"ur das Min-Max-Prinzip gilt
\begin{equation}
 \tilde x=\frac{x-x_{\min,\mathrm{train}}}
 {x_{\max,\mathrm{train}}-x_{\min,\mathrm{train}}}.
\end{equation}
Die Embedded-LSTM-Modelle verwenden dagegen die robuste Skalierung mit Trainingsmedian $m$ und Interquartilsabstand $s$, also $\tilde x=(x-m)/s$. Beide Verfahren lernen ihre Skalierungsparameter aus Trainingsdaten. Das Prinzip ist gemeinsam, das konkrete Verfahren muss je Modell erhalten bleiben.

\textbf{Pruning-Kriterium.} Die Bewertung eines rekurrenten Kanals $h$ fasst die Normen seiner Input- und rekurrenten Gewichtszeilen \"uber die vier LSTM-Bl\"ocke zusammen:
\begin{equation}
 s_h=\sum_{a\in\{i,f,o,g\}}
 \left(\left\|W_{\mathrm{ih}}^{(a)}[h,:]\right\|_2+
       \left\|W_{\mathrm{hh}}^{(a)}[h,:]\right\|_2\right).
\end{equation}
F\"ur die Menge $\mathcal R$ beibehaltener Kan\"ale werden neue dichte Matrizen gebildet:
\begin{align}
 \widetilde W_{\mathrm{ih}}^{(a)}&=W_{\mathrm{ih}}^{(a)}[\mathcal R,:],\\
 \widetilde W_{\mathrm{hh}}^{(a)}&=W_{\mathrm{hh}}^{(a)}[\mathcal R,\mathcal R],\\
 \widetilde W_{\mathrm{MLP},1}&=W_{\mathrm{MLP},1}[:,\mathcal R].
\end{align}
So wird aus dem allgemeinen Pruning-Bild eine konkrete, konsistente Verkleinerung der rekurrenten Implementierung.

\textbf{Gewichtsspeicherung.} F\"ur $N_W$ Gewichte mit $b$ Bits ist der reine Speicher $B_W=N_W b/8$ Byte. Im gemischten Export kommen FP32-Skalen, FP32-Teile und Programmcode hinzu. Darum ist die ideale Vierfachreduktion einzelner Matrizen nicht die Reduktion der gesamten Firmware.

\clearpage
\section*{Bildauswahl und Weg in die Folien}
\addcontentsline{toc}{section}{Bildauswahl und Folienplanung}
Die Originalabbildungen bleiben die Belege. F\"ur die sp\"atere PowerPoint werden bei dichten Mehrfachdiagrammen gezielt Teilpanels vergr\"o\ss ert. Es werden keine neuen Ergebniswerte oder vereinfachten Ersatzkurven erfunden. Neue Elemente dienen lediglich der Orientierung, zum Beispiel der Hervorhebung eines Pfads im BMS-Anforderungsbild.

\small
\begin{tabularx}{\linewidth}{@{}p{22mm}>{\raggedright\arraybackslash}X>{\raggedright\arraybackslash}p{63mm}@{}}
\toprule
Station & Originalabbildungen & Aufgabe im Vortrag \\
\midrule
1, 27 & 2.1 & Gemeinsamer Rahmen und R\"uckbindung. \\
Vor 4 & 2.1, Fokusansicht & Deployment und Modellkomplexit\"at, Accuracy erg\"anzend. \\
2, 3 & 4.1, 4.4, 4.6 & Datenvielfalt und reale BMS-Zielumgebung. \\
4 bis 7 & 3.2, 5.1, 5.2, 5.3 & Netzoperation, Merkmale, Historie und saubere Evaluation. \\
8 & 5.5 und 5.6 gemeinsam & Repr\"asentationswahl und erreichbare SOH-Sch\"atzung. \\
9 & 3.3 und 3.4 & Interne Zeitverarbeitung statt explizitem MLP-Historienvektor. \\
Vor 10 & 2.1, Fokusansicht & Performance und Hardware. \\
10, 11 & 6.1 und 6.2 & Modellfamilien und Testfamilien. \\
12 bis 15 & 6.4, 6.5, 6.9, 6.13 & Baseline, Stromfehler, Erholung und lokales Ereignis. \\
16, 17 & 6.15 und 6.16 & Einordnung aller Szenarien und Synthese. \\
18, 19 & 6.19 und zus\"atzlich 6.20 & Ressourcen und kontinuierliche Ausf\"uhrung. \\
Vor 20 & 2.1, Fokusansicht & Deployment und Hardware, Accuracy als Kontrollgr\"o\ss e. \\
20, 21 & 7.1, Daten aus 7.2, 7.3 & Konkrete Architektur und Optimierungspfade. \\
22, 23 & 3.5 und 3.7 & Technische Grundlagen direkt vor ihren Ergebnissen. \\
24 & 7.8 und zus\"atzlich 7.9 & SOC und SOH getrennt beurteilen. \\
25, 26 & 7.12, 7.13, Tabellen 7.1/7.2 & Speicher, Gesamtlatenz und reine Inferenz. \\
Reserve & 6.7 & Stromrauschen als erg\"anzender Mechanismus. \\
\bottomrule
\end{tabularx}
\normalsize

\textbf{Lesbarkeit.} Bei 4.4 nur den relevanten Zellverlauf und die Profilzuordnung fokussieren. Bei 5.5 und 5.6 dieselbe Seite beibehalten, aber die Heatmap und einen Ergebnisverlauf nacheinander hervorheben. Bei 6.5 die Fehlerzunahme und den SOC-Verlauf ausw\"ahlen. Bei 6.13 den Zeitverlauf und die ereignisbezogene Reaktion gegen\"uberstellen. Bei 7.8 und 7.9 jeweils zuerst den Boxplot zeigen. Das ganze Original dient als Kontext, ein vergr\"o\ss erter Ausschnitt macht den Befund lesbar.

\textbf{Navigation.} F\"ur die Folien den seitlichen Fortschrittsbalken aus v3 aufnehmen. Themen hei\ss en Anforderungen, Daten, zeitliche Repr\"asentation, Robustheit und Embedded-Optimierung. Keine Aufteilung in nummerierte Studien. Die drei Fokusansichten stehen vor MLP, Benchmark und Optimierung auf den Seiten \pageref{focus:mlp}, \pageref{focus:benchmark} und \pageref{focus:optimization}. Am Hardware-\"Ubergang innerhalb des Benchmarks wird der linke Bereich zus\"atzlich betont. Der Rest des Rahmens bleibt als Kontext sichtbar.

\textbf{Typografie.} Sp\"atere Folien mit Arial, Titeln in 24 pt und Text in 18 pt. Dichte Originaldiagramme nicht beliebig verkleinern. F\"ur sichtbare Diagrammtexte gelten dieselben Lesbarkeitsanforderungen. Kein dauerhafter roter Fazitsatz am Folienende.

\clearpage
\section*{Fachliche Pr\"ufhinweise zu den Originalen}
\addcontentsline{toc}{section}{Pruefhinweise zu den Originalen}
Diese Hinweise geh\"oren zur Vorbereitung, nicht in den gesprochenen Vortrag. Die Originaldateien bleiben unver\"andert. Die BMS-Fokusansichten erg\"anzen ausschlie\ss lich Konturen im Dokument. Der Abgleich hat folgende Punkte ergeben, die vor der finalen Folienproduktion relevant sind.

\textbf{Abbildung 7.12: RAM-Zahl nicht konsistent.} Der quantisierte SOC-Balken zeigt etwa 3,5 KB. Tabelle 7.1 nennt 3,96 KB. Deshalb wird diese einzelne RAM-Zahl in der Erz\"ahlung nicht als feststehender Wert wiederholt. Vor der \"Ubernahme des RAM-Panels m\"ussen Grafik und Tabelle mit dem zugeh\"origen Firmware-Ergebnis abgeglichen werden.

\textbf{Abbildung 7.12: Parameterzahl und Speicherbasis trennen.} Das obere rechte Panel ist eine idealisierte Gewichtsspeicherung und setzt beim INT8-Balken ungef\"ahr ein Viertel des gesamten FP32-Gewichtsspeichers an. Die real untersuchte Firmware quantisiert dagegen nur die rekurrenten Matrizen. Das untere linke Panel zeigt ihren Flash-Bedarf. Die oberen Parameterzahlen sind zudem nicht ohne Weiteres aus den aktuellen Dimensionen in Abschnitt 7.3 rekonstruierbar. Sie sollten vor einer Formelherleitung am konkreten Export gepr\"uft werden. Die Erz\"ahlung verwendet die im Methodenteil explizit angegebenen MAC-Zahlen und keine aus dem Balken abgelesene Parameterherleitung.

\textbf{Abbildung 7.8: kleine Zahlenabweichung.} Die Originallegende nennt 2,33 und 2,78 Prozentpunkte f\"ur Pruned und Quantized. Der laufende Dissertationstext nennt 2,34 und 2,79. Die Erz\"ahlung verwendet deshalb f\"ur diesen Vergleich transparent gerundete Werte von 2,3 und 2,8. Die Abweichung von jeweils 0,01 Prozentpunkten bleibt ein abzugleichender Quellenpunkt, kein nachgewiesener anderer Effekt.

\textbf{Abbildung 7.13: Messgrenze korrekt.} Gezeigt wird die Host-Gesamtlatenz einschlie\ss lich UART. Die Kernelzeiten stammen aus Tabellen 7.1 und 7.2. Beide werden im Vortrag ausdr\"ucklich getrennt. Die Zeitwerte charakterisieren die verwendeten Referenzkernels und Compilerbedingungen, nicht eine allgemeine Grenze f\"ur INT8-Inferenz auf Mikrocontrollern.

\textbf{Korrelation und Robustheitsmechanismen.} 5.1 zeigt lineare Korrelationen mit kumulierten Ladungsmengen, nicht eine kausale Stromwirkung oder gemessene Feature Importance. 6.5 zeigt Gain, nicht den additiven Offset. 6.9 betrifft definierte Erholung relativ zum ungest\"orten Vergleich, nicht garantierte absolute SOC-Genauigkeit. 6.13 zeigt einen Spannungsimpuls. Stromintegration und geringe Spannungsempfindlichkeit haben deshalb unterschiedliche Erkl\"arungen. Der breitere Kontext des DD ist eine begr\"undete strukturelle Interpretation, kein separater Ablationsnachweis.

\textbf{Verschiedene Untersuchungen nicht gleichsetzen.} Kapitel 6 verwendet f\"ur die Robustheit einen GRU-SOC-Zweig im Fensterbetrieb und einen gemeinsamen st\"undlichen SOH-LSTM. Kapitel 7 untersucht andere, schrittweise arbeitende LSTM-SOC- und LSTM-SOH-Modelle. Ihre MAE-Werte liefern keinen unmittelbaren Architekturvergleich. Die eigene H755-BMS-Platine und die H753-Hardwarebenchmarks sind ebenfalls getrennte Teile der experimentellen Grundlage.

\textbf{Scalability richtig einordnen.} Eine Architektur an unterschiedliche Ressourcenbudgets anzupassen, ist ein Deployment-Aspekt. Pruning kann das unterst\"utzen. Eine nachgewiesene \"Ubertragung auf beliebige Controller, Zellchemien oder gr\"o\ss ere Packs folgt daraus jedoch nicht. Deshalb werden diese nicht als bereits experimentell belegte Skalierungsergebnisse dargestellt.

\textbf{Umfang dieser Bearbeitung.} Dissertation, Originalbilder und vorhandene PowerPoints wurden nicht ge\"andert. Die vorige Erz\"ahlfassung v2 bleibt im Sicherungsordner erhalten:
\par\texttt{\_build\_erzaehlfassung/previous\_v2}

%TC:endignore
\end{document}
